% Shared preamble for the supporting-information documents (S1-S5).
\documentclass[10pt,letterpaper]{article}
\usepackage[margin=1in]{geometry}
\usepackage{amsmath,amssymb}
\usepackage{booktabs,array,longtable}
\usepackage[table]{xcolor}
\usepackage{hyperref}
\usepackage{microtype}
\usepackage{listings}
\usepackage{graphicx}
\usepackage[labelfont=bf,labelsep=period,justification=raggedright,singlelinecheck=off]{caption}
\definecolor{linkblue}{HTML}{1C5CAB}
\definecolor{codebg}{HTML}{F4F4F1}
\hypersetup{colorlinks=true, linkcolor=linkblue, citecolor=linkblue, urlcolor=linkblue}
\lstset{basicstyle=\ttfamily\footnotesize, backgroundcolor=\color{codebg},
  frame=none, breaklines=true, columns=fullflexible, keepspaces=true,
  xleftmargin=4pt, xrightmargin=4pt, aboveskip=6pt, belowskip=6pt}
\setlength{\parindent}{0pt}
\setlength{\parskip}{5pt}
\raggedright
\renewcommand{\tablename}{Table}
\newcommand{\sititle}[2]{{\Large\bfseries #1}\par\vspace{3pt}{\large #2}\par\vspace{2pt}
  {\small\color{black!60}Supporting information for ``Environmental fluctuation limits gene regulation unless it can be tracked''}\par\vspace{10pt}}

\begin{document}
\sititle{S4 Text.}{Earlier EvoSim results and the audit trail}

This supplement records the EvoSim results that preceded the present study, what
became of each, and the measurement errors found along the way together with the
check that exposed each one. The complete record, including every retracted
analysis, is in the project's \texttt{README.md} and \texttt{PAPER.md}.

\section*{1. Earlier EvoSim results and their status}
These results came from experiments with smaller populations (mostly between 2{,}000 and 3{,}000 organisms) than the 180-run batch in the main text.

\begin{longtable}{@{}p{0.27\linewidth}p{0.36\linewidth}p{0.3\linewidth}@{}}
\toprule
Earlier result & Evidence & Status after this study\\
\midrule
\endhead
Fluctuation increases the number of network nodes &
7 of 7 frozen--fluctuating comparisons positive and significant; mean $+6.3$
nodes (range $+2.6$ to $+10.2$) &
Reproduces published results~[1,2]. R-Aevol shows the opposite (P4), at every
grain tested (G1--G4), so this is a difference between the models.\\
Fluctuation increases connection density &
Count density $+0.09$ to $+0.14$, significant in 5 of 6 comparisons; the edge
increase exceeded that expected from node number alone by 27 to 41 percentage
points; stable across development iterations 4 to 20 &
The count effect replicates ($+0.12$ at grid 128), but strength density falls
by 11\% and total strength does not change. The interpretation as denser wiring
is withdrawn.\\
The information content of the cue does not matter &
Pooled over 15 informative--scrambled comparisons, mean difference $-0.17$ nodes,
$t=-0.30$, sign test $p=0.61$ &
Holds in EvoSim, where organisms exploit about 11\% of the cue's information.
Contradicted in R-Aevol, where the cue is used (P2).\\
Wiring cost makes fluctuation lethal &
In the sparsest regime tested, 0 of 60 frozen and 37 and 33 of 60 fluctuating
populations went extinct (Fisher $p=3.6\times10^{-13}$) &
Consistent with the viability pilot for P5: fluctuating populations survived 2,
2 and 5 of 6 times at grids 80, 96 and 112, against 6 of 6 frozen.\\
Task difficulty does not set genome size &
Four manipulations of the task, all null (detection floors 2.2 to 4.2 nodes) &
Stands.\\
Organisms read the cue when development is reversible &
With living organisms re-developed every few steps, switchers 0.65 against 0.23
(informative against scrambled, 100 populations per arm); genome size unchanged
($p=0.84$, 95\% CI $-2.1$ to $+2.6$ nodes) &
Stands. The strength measure has not been taken in this regime.\\
Long-run outcome &
8 runs of 25 million steps (about 630{,}000 generations): node number reached the
minimum (12 at the time) by generation 14{,}000 and stayed there &
Stands; the present version has a minimum of 13 because a sensor node was
added.\\
\bottomrule
\end{longtable}

\section*{2. Record of pre-registered predictions}
Throughout the project, predictions were committed to version control before the
runs that tested them.

\begin{center}
\small
\begin{tabular}{@{}p{0.62\linewidth}p{0.3\linewidth}@{}}
\toprule
Prediction & Outcome\\
\midrule
\multicolumn{2}{@{}l}{\emph{Earlier EvoSim work}}\\
Genome size rises with the complexity of the task & wrong (no change)\\
Rotating the yields alone makes the task conditional & wrong (inflow dominated)\\
An abiotic cue will be read because it cannot be depleted & wrong (legible but unread)\\
A conditional genome is adaptive but hard to reach by mutation & wrong (selected against)\\
Reversible development makes a fine-grained regime habitable & correct (survival rose from 0 of 3 to 3 of 3)\\
A conditional response, once evolved, raises genome size & wrong (the response evolved, size did not change)\\
\multicolumn{2}{@{}l}{\emph{This study}}\\
P1--P4 (R-Aevol sweep) & all supported\\
P5 (EvoSim confirmation) & not run (no viable regime)\\
G1, G3 (grain, coarse against fine) & inconclusive\\
G2, G4 (grain, coarse against frozen) & refuted\\
\bottomrule
\end{tabular}
\end{center}

\section*{3. Measurement errors and the checks that found them}
Almost every wrong conclusion in the project traced back to an error of measurement, and each was found by a check whose correct answer was known in advance. They are listed because the same errors recur in simulation
studies.

\begin{longtable}{@{}p{0.3\linewidth}p{0.34\linewidth}p{0.28\linewidth}@{}}
\toprule
Error & What it produced & How it was found\\
\midrule
\endhead
Edge counts used as a measure of wiring & A density increase that is not an
increase in wiring strength (this study) & Measuring summed absolute weight,
prompted by the R-Aevol replication\\
A statistic read without checking its definition & \texttt{mean\_abs\_w} is summed
weight divided by all $n_g^2$ slots, so multiplying it by the edge count gave
meaningless ``total strength'' columns & Reading the source of the statistic
before deriving from it\\
Parameters attributed to the wrong source & A tenfold error in the stated
mutation rate of R-Aevol's configuration, in a draft of the Methods & Comparing
against the platform's example files\\
A control that was not a control & Turning off the yield rotation also made the world poorer, and the arm with the least conditional task had the largest genomes & Measuring population size and energy
in every arm\\
A no-op manipulation that was not a no-op & Removing a node's inputs changed its
output downstream and altered every trait (repair 0.995 to 0.127) & Requiring a
zero-effect arm whose result is known\\
Tags that leaked & Tagging organisms by storage slot, which is reused after
death, spread the tag to unrelated organisms & A control arm that reached
fixation when it should not\\
Wrong unit of analysis & Pooling repeated measurements of the same node inflated
a correlation & Reducing each node to one value\\
Survivorship in the sampling frame & Sampling only at the end of runs, when almost
all nodes were old & Sampling during runs\\
A derived timescale from the wrong statistic & Generations per environmental
state computed from a survivor-biased lifespan, a threefold error & Using the
slope of lineage depth\\
Small samples & Three promising signals from 3 or 4 seeds, all gone at 12 seeds &
Replication at larger $n$\\
A control checked in one mode only & The scrambled cue leaked 0.067 bits of
information without a spatial gradient, the mode mostly used & Measuring its
information content in every mode; fixed to 0.0004 bits\\
A regime that kills one arm & At grid 48, 6 of 6 fluctuating and 0 of 6 frozen
populations died, which would have made the comparison one of survivors & A
viability check at full run length before any batch\\
Identical arms counted as replicates & With a single state the cue is zero, so
the three frozen arms were the same runs & Checking that arms differ\\
A verification script that could not fail & A scoping error kept an accumulator
at its initial value, so the script would have reported success whatever the
data & A negative control with a known failure\\
A significance test against the wrong spread & Dividing a difference by one arm's standard deviation instead of the standard error of the difference reported ``no effect'' for $t=4.2$ & Recomputing with Welch and permutation tests\\
\bottomrule
\end{longtable}

\section*{References}
\small
1. Tsuda ME, Kawata M. Evolution of gene regulatory networks by fluctuating
selection and intrinsic constraints. PLoS Comput Biol. 2010;6(8):e1000873.\\
2. Bentkowski P, van Oosterhout C, Mock T. A model of genome size evolution for
prokaryotes in stable and fluctuating environments. Genome Biol Evol.
2015;7(8):2344--51.

\end{document}
